% =============================================================================
%  main.tex — Exemplo de apresentação (template IComp não oficial)
%  Compile: pdflatex main.tex (2x)
% =============================================================================
\documentclass[aspectratio=169, 10pt, brazil]{beamer}

\usepackage[T1]{fontenc}
\usepackage[utf8]{inputenc}
\usepackage[brazil]{babel}
\usepackage{beamerthemeIComp}
\usepackage{booktabs, multirow, colortbl}
\usepackage{amsmath, amssymb}
\usepackage{tikz}
\usetikzlibrary{arrows.meta, positioning, shapes.geometric}

% ---------------------------------------------------------------------------
%  METADADOS — edite aqui
% ---------------------------------------------------------------------------
\title{Adaptação de Domínio Livre de Fonte em\\
       Reconhecimento de Ações em Vídeo\\
       Baseada em Protótipos Multicêntricos}
\subtitle{Exame de Qualificação — Mestrado em Informática}
\author[M. Vanderlei]{Miguel Vanderlei de Oliveira}
\institute[IComp · PPGI · UFAM]{%
  Instituto de Computação — IComp\\
  Programa de Pós-Graduação em Informática — PPGI · UFAM
}
\orientador{Orientadora: Prof.\textsuperscript{a} Dr.\textsuperscript{a} Eulanda Miranda dos Santos}
\contato{miguel.vanderlei@icomp.ufam.edu.br}
\rodape{IComp · PPGI · UFAM}
\date{Manaus, Julho de 2026}

% ---------------------------------------------------------------------------
\begin{document}

% ── Capa ────────────────────────────────────────────────────────────────────
\begin{frame}[noframenumbering, plain]
  \titlepage
\end{frame}

% ── Sumário ─────────────────────────────────────────────────────────────────
\begin{frame}{Sumário}
  \begin{columns}[T,onlytextwidth]
    \begin{column}{0.47\textwidth}
      \tableofcontents[sections={1-3},hideallsubsections]
    \end{column}
    \begin{column}{0.47\textwidth}
      \tableofcontents[sections={4-6},hideallsubsections]
    \end{column}
  \end{columns}
\end{frame}

% ============================================================================
\section{Introdução e Motivação}
% ============================================================================

\begin{frame}{Reconhecimento de Ações em Vídeo}
  \begin{columns}[T]
    \begin{column}{0.58\textwidth}
      \begin{itemize}
        \item Tarefa fundamental da Visão Computacional
        \item Backbone padrão: \hgreen{I3D} (Inflated 3D ConvNet)
        \item Excelente acurácia em UCF101, HMDB51 etc.
      \end{itemize}
      \vspace{0.8em}
      \begin{block}{Problema — \textit{Domain Shift}}
        Queda de desempenho ao mudar câmera, iluminação
        ou estilo de execução entre domínios.
      \end{block}
    \end{column}
    \begin{column}{0.38\textwidth}
      \centering
      \begin{tikzpicture}[node distance=0.9cm]
        \node[draw=ufamNavy, fill=ufamLightGray, rounded corners=3pt,
              align=center, font=\small\bfseries, text=ufamNavy,
              minimum width=2.8cm, minimum height=0.9cm] (src)
          {Fonte\\(UCF101)};
        \node[draw=ufamGreen, fill=ufamGreenLight, rounded corners=3pt,
              align=center, font=\small\bfseries, text=ufamGreenDark,
              minimum width=2.8cm, minimum height=0.9cm, below=of src] (tgt)
          {Alvo\\(HMDB51)};
        \draw[-{Stealth}, line width=1.2pt, ufamRed, dashed]
          (src) -- node[right, font=\tiny, text=ufamRed]
          {$\downarrow$ acc.} (tgt);
      \end{tikzpicture}
    \end{column}
  \end{columns}
\end{frame}

\begin{frame}{Do UDA ao SFDA}
  \begin{columns}[T]
    \begin{column}{0.47\textwidth}
      \begin{alertblock}{UDA — Limitação}
        Exige dados da fonte durante a adaptação.\\
        Inviável por \textbf{privacidade} e \textbf{armazenamento}.
      \end{alertblock}
    \end{column}
    \begin{column}{0.47\textwidth}
      \begin{exampleblock}{SFDA — Solução}
        (1) Pré-treinamento supervisionado na \textbf{fonte}\\
        (2) Adaptação no \textbf{alvo} sem dados da fonte
      \end{exampleblock}
    \end{column}
  \end{columns}
  \vspace{-0.2em}
  \begin{block}{Estado da Arte — \textbf{CleanAdapt} (Dasgupta et al., 2025)}
    Trata adaptação como aprendizado com rótulos ruidosos.
    Seleciona fração $\tau$ de menor perda por classe (\emph{small-loss trick}).\\[0.3em]
    \hred{Limitação:} trade-off \textbf{qualidade} $\times$ \textbf{cobertura} —
    penaliza classes minoritárias e cauda longa.
  \end{block}
\end{frame}

\begin{frame}{Objetivos}
  \begin{block}{Objetivo Geral}
    Propor método de SFVDA com \hgreen{protótipos multicêntricos} que supere
    a limitação de cobertura do CleanAdapt via \textbf{re-rotulagem geométrica}
    no espaço latente.
  \end{block}
  \vspace{0.5em}
  \textbf{Objetivos Específicos:}
  \begin{enumerate}
    \item Analisar o trade-off qualidade $\times$ cobertura do CleanAdapt
    \item Propor agrupamento multicêntrico com re-amostragem balanceada
    \item Implementar re-rotulagem por votação majoritária (KNN + K-Means)
    \item Investigar estratégias \textit{imbalance-aware} para cauda longa
  \end{enumerate}
\end{frame}

% ============================================================================
\section{Fundamentos}
% ============================================================================

\begin{frame}{Fundamentos — Visão Geral}
  \begin{columns}[T]
    \begin{column}{0.47\textwidth}
      \begin{enumerate}
        \item \textbf{Aprendizado Profundo / CNN}\\
              \small Entropia cruzada; SGD; overfitting
        \item \textbf{Arquitetura I3D}\\
              \small Inflated 3D ConvNet; pré-treino Kinetics
        \item \textbf{SFDA / SFVDA}\\
              \small Dois estágios; sem dados fonte
      \end{enumerate}
    \end{column}
    \begin{column}{0.47\textwidth}
      \begin{enumerate}\setcounter{enumi}{3}
        \item \textbf{Pseudo-rótulos e Small-Loss}\\
              \small Efeito memorização (Arpit et al., 2017)
        \item \textbf{Estrutura Prof.–Aluno + EMA}\\
              \small $\theta_T \leftarrow \alpha\theta_T + (1-\alpha)\theta_S$
        \item \textbf{K-Means e Protótipos}\\
              \small Agrupamento latente; viés de densidade
      \end{enumerate}
    \end{column}
  \end{columns}
\end{frame}

% ============================================================================
\section{Trabalhos Relacionados}
% ============================================================================

\begin{frame}{Posicionamento na Literatura}
  \small\renewcommand{\arraystretch}{1.3}
  \begin{tabular}{@{}llp{3.5cm}@{}}
    \toprule
    \rowcolor{ufamNavy}
    \color{white}\textbf{Método} &
    \color{white}\textbf{Estratégia} &
    \color{white}\textbf{Limitação} \\
    \midrule
    SHOT          & Centroides + inf.\ mútua   & Só imagens \\
    \rowcolor{ufamLightGray}
    ATCON         & Consistência temporal      & Rótulos globais inalterados \\
    STHC          & Banco de memória histórico & Sem topologia de classes \\
    \rowcolor{ufamLightGray}
    BMD           & Agrupamento balanceado     & Estima distrib.\ do alvo \\
    DCPL          & Matriz de transição ruído  & Rótulos fixos \\
    \rowcolor{ufamLightGray}
    \textbf{CleanAdapt} & \textbf{Small-loss + Prof.–Aluno} &
      \textbf{Trade-off qualidade $\times$ qtd.} \\
    \rowcolor{ufamGreenLight}
    \textbf{Proposta} & \textbf{Protótipos multicêntricos} &
      Agrupamento em cauda longa \\
    \bottomrule
  \end{tabular}
\end{frame}

% ============================================================================
\section{Abordagem Proposta}
% ============================================================================

\begin{frame}{Pipeline — Visão Geral}
  \centering
  \begin{tikzpicture}[
      node distance=0.65cm and 1.0cm,
      box/.style={draw=ufamNavy, rounded corners=3pt, fill=ufamLightGray,
                  minimum width=2.0cm, minimum height=0.65cm,
                  align=center, font=\scriptsize},
      gbox/.style={box, fill=ufamGreenLight, draw=ufamGreenDark},
      tbox/.style={box, fill=ufamTealLight, draw=ufamTeal},
      arr/.style={-{Stealth[length=4.5pt]}, line width=0.9pt, ufamNavy},
      garr/.style={arr, ufamGreen},
    ]
    \node[tbox] (src)  {Fonte I3D};
    \node[box, right=of src]  (prof) {Professor\\(aug.\ fraco)};
    \node[box, below=0.5cm of prof] (aluno) {Aluno\\(aug.\ forte)};
    \node[gbox, right=of prof] (dcl)  {$\mathcal{D}_{cl}$\\confiável};
    \node[gbox, right=of aluno] (dun) {$\mathcal{D}_{un}$\\incerto};
    \node[tbox, right=of dcl]  (proto) {Protótipos\\K-Means};
    \node[gbox, right=of dun]  (relabel) {Re-rotulagem\\KNN};
    \node[box, right=1.2cm of proto] (train) {Treino\\$\mathcal{D}_{cl}\cup\mathcal{D}_{un}^{\text{rel}}$};
    \draw[arr] (src)   -- (prof);
    \draw[arr] (prof)  -- node[above,font=\tiny]{$\tau$} (dcl);
    \draw[arr] (prof)  -- node[below,font=\tiny]{alta perda} (dun);
    \draw[arr] (dcl)   -- (proto);
    \draw[arr] (proto) -- (relabel);
    \draw[arr] (dun)   -- (relabel);
    \draw[arr] (dcl.east) to[bend left=18] (train.north);
    \draw[arr] (relabel)  -- (train);
    \draw[garr] (train.south) to[bend right=25]
      node[below,font=\tiny]{backprop} (aluno.south);
    \draw[garr] (aluno.west) to[bend right=15]
      node[left,font=\tiny]{EMA} (prof.west);
  \end{tikzpicture}
\end{frame}

\begin{frame}{Os Quatro Passos da Re-rotulagem}
  \begin{enumerate}
    \item \textbf{Subconjunto confiável} — Professor filtra fração $\tau$ de
          menor perda por classe $\Rightarrow$ $\mathcal{D}_{cl}$.
          Complemento: $\mathcal{D}_{un} = \mathcal{D}_t \setminus \mathcal{D}_{cl}$.

    \item \textbf{Protótipos multicêntricos} — K-Means com $K$ centroides sobre
          embeddings $\ell_2$-normalizados de $\mathcal{D}_{cl}^c$ para cada
          classe $c \in \mathcal{Y}$.

    \item \textbf{Conjunto global} $\mathcal{P} = \bigcup_c \mathcal{P}_c$ —
          métrica \textbf{Mahalanobis} (selecionada via Optuna).

    \item \textbf{Re-rotulagem por votação}
          $\hat{y}_i^* = \arg\max_c |\{\mu \in \mathcal{N}_k(x_i,\mathcal{P})
          : \text{classe}(\mu)=c\}|$

          Aluno treina sobre $\mathcal{D}_{cl} \cup \mathcal{D}_{un}^{\text{rel}}$.
  \end{enumerate}
  \vspace{0.4em}
  \begin{exampleblock}{Desacoplamento chave}
    $\tau$ controla apenas o \emph{tamanho do âncora}, não a cobertura do treino.
  \end{exampleblock}
\end{frame}

\begin{frame}{Robustez — Distribuições de Cauda Longa}
  \begin{columns}[T]
    \begin{column}{0.48\textwidth}
      \begin{alertblock}{Problema}
        K-Means tem \textbf{viés de densidade}: classes majoritárias
        dominam os centroides.
      \end{alertblock}
      \vspace{0.5em}
      \begin{exampleblock}{Re-amostragem Balanceada}
        $n_{\min} = \min_{c} |\mathcal{D}_{cl}^c|$ instâncias por classe
        $\Rightarrow$ todas as classes contribuem igualmente para $\mathcal{P}$.
      \end{exampleblock}
    \end{column}
    \begin{column}{0.48\textwidth}
      \begin{block}{Direções Futuras}
        \begin{itemize}
          \item Amostragem curricular — ISFDA
          \item Interseção de vizinhanças — Robust NN
          \item Classificador normalizado $\ell_2$ — ICPR
        \end{itemize}
      \end{block}
    \end{column}
  \end{columns}
\end{frame}

% ============================================================================
\section{Resultados Preliminares}
% ============================================================================

\begin{frame}{Resultados — UCF101 $\leftrightarrow$ HMDB51}
  \centering\small\renewcommand{\arraystretch}{1.15}
  \begin{tabular}{@{}lcc@{}}
    \toprule
    \rowcolor{ufamNavy}
    \color{white}\textbf{Método} &
    \color{white}\textbf{UCF$\!\to\!$HMDB (\%)} &
    \color{white}\textbf{HMDB$\!\to\!$UCF (\%)} \\
    \midrule
    CleanAdapt+TS (baseline) & 87{,}78 & 96{,}67 \\
    \rowcolor{ufamLightGray}
    Aumento linear de $\tau$  & 87{,}50 & 96{,}50 \\
    ATCON$^\dagger$           & 79{,}70 & 85{,}30 \\
    \rowcolor{ufamLightGray}
    MTRAN$^\ddagger$          & 92{,}20 & 95{,}30 \\
    \midrule
    \rowcolor{ufamGreenLight}
    \textbf{Protótipos (proposta)} &
      \textbf{92{,}78} & \textbf{98{,}25} \\
    \bottomrule
    \multicolumn{3}{l}{\tiny $^\dagger$ Backbone TRN;\ \
    $^\ddagger$ Transformer multimodal}
  \end{tabular}
  \vspace{0.2em}
  \begin{columns}[T]
    \begin{column}{0.47\textwidth}
      \begin{exampleblock}{UCF $\to$ HMDB}
        \textbf{+5,00 p.p.}\ sobre baseline\\
        (direção mais difícil; alvo desbalanceado)
      \end{exampleblock}
    \end{column}
    \begin{column}{0.47\textwidth}
      \begin{exampleblock}{HMDB $\to$ UCF}
        \textbf{+1,58 p.p.}\ sobre baseline\\
        Competitivo com Transformers via I3D
      \end{exampleblock}
    \end{column}
  \end{columns}
\end{frame}

\begin{frame}{Perguntas de Pesquisa}
  \begin{columns}[T,onlytextwidth]
    \begin{column}{0.47\textwidth}
      \begin{block}{PP1 — Sensibilidade ao $\tau$}
        \begin{itemize}
          \item Proposta: $\tau$ regula âncora, não cobertura
          \item Melhor UCF$\!\to\!$HMDB com $\tau\!=\!0{,}2$
          \item Âncora pequeno + re-rotulagem = suficiente
        \end{itemize}
      \end{block}
    \end{column}
    \begin{column}{0.47\textwidth}
      \begin{block}{PP2 — K$>$1 vs K$=$1}
        \begin{itemize}
          \item K$>$1 supera K$=$1 em até +1,94 p.p.\ em $\mathcal{D}_{un}$
          \item Múltiplos centroides capturam sub-modos intra-classe
        \end{itemize}
      \end{block}
    \end{column}
  \end{columns}
\end{frame}

\begin{frame}{Epic-Kitchens}{Limitação em distribuições de cauda longa}
  \begin{columns}[T,onlytextwidth]
    \begin{column}{0.47\textwidth}
      \begin{alertblock}{Resultado}
        Baseline supera proposta nos 6 cenários, com lacuna de até
        \textbf{16 p.p.}
      \end{alertblock}
    \end{column}
    \begin{column}{0.47\textwidth}
      \begin{block}{Diagnóstico}
        Viés de densidade do K-Means e pseudo-rótulos iniciais ruidosos.
      \end{block}
      \begin{exampleblock}{Direção}
        Estratégias \textit{imbalance-aware} nos próximos passos.
      \end{exampleblock}
    \end{column}
  \end{columns}
\end{frame}

% ============================================================================
\section{Considerações Parciais e Próximos Passos}
% ============================================================================

\begin{frame}{Contribuições e Cronograma}
  \begin{columns}[T]
    \begin{column}{0.47\textwidth}
      \textbf{Contribuições:}
      \begin{itemize}
        \item Re-rotulagem geométrica multicêntrica
        \item +5 p.p.\ sobre SOTA em UCF$\leftrightarrow$HMDB
        \item Desacoplamento âncora/cobertura via $\tau$
        \item \hgreen{Artigo aceito no BRACIS 2026}
      \end{itemize}
    \end{column}
    \begin{column}{0.47\textwidth}
      \textbf{Próximos Passos:}
      \begin{enumerate}
        \item Validar re-amostragem no Epic-Kitchens
        \item Amostragem curricular (ISFDA)
        \item Interseção de vizinhanças (Robust NN)
        \item Comparação ampliada + redação final (2027)
      \end{enumerate}
    \end{column}
  \end{columns}
\end{frame}

% ── Slide Final ──────────────────────────────────────────────────────────────
\textofinal{Obrigado!}
\apoio{%
  \imgapoio[height=0.55cm]{logo_cnpq.png}
  \imgapoio[height=0.55cm]{logo_capes.png}
  \imgapoio[height=0.72cm]{logo_fapeam.png}
}
\agradecimentos

\end{document}
